\paragraph{Setup.}  The head-to-head runs used Google Colab machines with one
physical core (two hyperthreads) and 12\,GB of memory.  On every machine the
two solvers ran one after the other, each pinned to the same logical CPU while
the rest of the machine was idle, with the same wall-clock budget of
\SI{600}{s}.  KaPoCE was stopped as in PACE~2021, by SIGTERM at the limit; it
normally stops itself after 590\,s.  Its seed, fixed in the released code, was
set to the run seed (a one-line change), so that its ten runs per graph are
independent.  Its output was checked: an edit set is valid iff the edited graph
is a disjoint union of cliques.  A run in which KaPoCE returned no valid
solution would count as a win for PXMem, as in the PACE rules; this did not
occur in any of the 150 runs.  PXMem ran with the released configuration.  The
Numba kernels were compiled once per machine before the timed runs, as KaPoCE
is compiled ahead of time.  We used all fifteen SNAP graphs, ten seeds each.
The configuration of PXMem was developed on eleven of these graphs; the PACE
heuristic-track instances below were not used during development.  We report
the mean paired difference $\bar\Delta$ (PXMem minus KaPoCE, in
disagreements) with a 95\% bootstrap interval, wins, ties and losses, the
two-sided sign test and its Holm correction over the fifteen graphs.

\begin{table}[t]
\centering
\caption{PXMem against KaPoCE, \SI{600}{s} per run, run one after the other on
the same machine and CPU (ten seeds per graph).  Mean costs; $\bar\Delta$: mean
paired difference PXMem$-$KaPoCE with a 95\% bootstrap interval; W/T/L: wins,
ties and losses of PXMem; $p$: two-sided sign test; $p_{\mathrm{Holm}}$: Holm
correction over the fifteen graphs (bold if below $0.05$); $t_{=}$: median time
at which PXMem first reached the final cost of KaPoCE (``--'' if it did not in
most runs).  KaPoCE returned a valid solution in every run.}\label{tab:h2h}
\small
\begin{tabular}{lrrrrrrrr}
\toprule
graph & $n$ & PXMem & KaPoCE & $\bar\Delta$ [95\% CI] & W/T/L & $p$ & $p_{\mathrm{Holm}}$ & $t_{=}$ (s) \\
\midrule
\texttt{ca-GrQc} & \num{5242} & \num{6050.0} & \num{6050.0} & $+0.0$  & 0/10/0 & 1.00 & 1.00 & 69 \\
Bitcoin-Alpha$^+$ & \num{3683} & \num{11437.0} & \num{11437.0} & $+0.0$  & 0/10/0 & 1.00 & 1.00 & 71 \\
Bitcoin-OTC$^+$ & \num{5573} & \num{16462.0} & \num{16462.0} & $+0.0$  & 0/10/0 & 1.00 & 1.00 & 181 \\
\texttt{ca-HepTh} & \num{9877} & \num{15813.0} & \num{15813.5} & $-0.5$ $[-0.8, -0.2]$ & 5/5/0 & 0.06 & 0.62 & 125 \\
\texttt{ca-HepPh} & \num{12008} & \num{46133.0} & \num{46133.0} & $+0.0$  & 0/10/0 & 1.00 & 1.00 & 107 \\
\texttt{ca-AstroPh} & \num{18772} & \num{125543.9} & \num{125606.8} & $-62.9$ $[-78.1, -49.5]$ & 10/0/0 & \textbf{0.002} & \textbf{0.029} & 62 \\
\texttt{ca-CondMat} & \num{23133} & \num{56834.3} & \num{56845.6} & $-11.3$ $[-16.4, -6.0]$ & 8/2/0 & \textbf{0.008} & 0.10 & 56 \\
\texttt{email-Enron} & \num{36692} & \num{150428.2} & \num{150460.8} & $-32.6$ $[-47.5, -17.1]$ & 9/0/1 & \textbf{0.021} & 0.26 & 139 \\
\texttt{loc-Brightkite} & \num{58228} & \num{174700.1} & \num{174703.2} & $-3.1$ $[-6.0, -0.4]$ & 7/0/3 & 0.34 & 1.00 & 432 \\
\texttt{Slashdot+} & \num{75144} & \num{351865.0} & \num{351869.1} & $-4.1$ $[-8.7, +0.1]$ & 6/2/2 & 0.29 & 1.00 & 579 \\
\texttt{soc-Epinions} & \num{75879} & \num{365164.5} & \num{365167.6} & $-3.1$ $[-22.7, +16.1]$ & 4/0/6 & 0.75 & 1.00 & -- \\
\texttt{Epinions+} & \num{114467} & \num{521817.3} & \num{521818.1} & $-0.8$ $[-17.6, +16.6]$ & 5/0/5 & 1.00 & 1.00 & -- \\
\texttt{com-DBLP} & \num{317080} & \num{606753.1} & \num{606787.2} & $-34.1$ $[-56.5, -15.1]$ & 9/0/1 & \textbf{0.021} & 0.26 & 296 \\
\texttt{com-Amazon} & \num{334863} & \num{612328.7} & \num{612311.6} & $+17.1$ $[+14.5, +20.0]$ & 0/0/10 & \textbf{0.002} & \textbf{0.029} & -- \\
\texttt{com-Youtube} & \num{1134890} & \num{2650090.2} & \num{2649791.7} & $+298.5$ $[+40.0, +574.9]$ & 4/0/6 & 0.75 & 1.00 & -- \\
\bottomrule
\end{tabular}

\end{table}

\paragraph{Solution quality at \SI{600}{s}.}  Table~\ref{tab:h2h} shows that on
the four smallest graphs and on \texttt{ca-HepPh} both solvers find the same
solutions in every run; PXMem reaches them after one to three minutes.  PXMem
is better in every run on \texttt{ca-AstroPh} (by 63 disagreements on average)
and in eight to nine of ten runs on \texttt{ca-CondMat}, \texttt{email-Enron}
and \texttt{com-DBLP} (by 11 to 34).  These four differences are significant by
the sign test, but after the Holm correction over fifteen graphs only
\texttt{ca-AstroPh} remains so.  On the four social networks of medium size the
two solvers are statistically indistinguishable.  KaPoCE is better in every run
on \texttt{com-Amazon} (by 17 disagreements, $0.003\%$), which survives the
correction, and on average on \texttt{com-Youtube}, where the runs of both
solvers vary by several hundred disagreements and the sign test is not
significant.  Over all 150 runs PXMem wins 67, ties 49 and loses 34.  The
differences are small in relative terms (at most $0.05\%$), far below the
certified gaps of Section~\ref{sec:snap}; they matter only when the last
fraction of a percent is wanted.  On \texttt{com-Amazon} we compared the two solutions of a local \SI{300}{s}
run block by block with Theorem~\ref{thm:px}.  They differed in more than $16{,}000$ blocks, almost all
ties, and KaPoCE was better in $24$ blocks of $8$ to $77$ vertices, each
containing a vertex of degree $22$ to $142$.  Solving these blocks exactly
with the MIP of Section~\ref{sec:gap} recovered the value of KaPoCE in $23$ of
them, mostly in less than a second, so the difficulty lies in finding the blocks,
not in solving them.

\begin{table}[t]
\centering
\caption{Mean relative difference (\%) of PXMem against KaPoCE for different
budgets (two runs per graph at \SI{60}{s} and \SI{150}{s}, ten at \SI{600}{s});
negative values favour PXMem.}\label{tab:budget}
\small
\begin{tabular}{lrrrrrr}
\toprule
 & \multicolumn{2}{c}{\SI{60}{s}} & \multicolumn{2}{c}{\SI{150}{s}} & \multicolumn{2}{c}{\SI{600}{s}} \\
\cmidrule(lr){2-3}\cmidrule(lr){4-5}\cmidrule(lr){6-7}
graph & med.\ (\%) & W/T/L & med.\ (\%) & W/T/L & med.\ (\%) & W/T/L \\
\midrule
\texttt{ca-GrQc} & $+0.000$ & 0/10/0 & $+0.000$ & 0/10/0 & $+0.000$ & 0/10/0 \\
Bitcoin-Alpha$^+$ & $+0.000$ & 2/6/2 & $+0.000$ & 0/8/2 & $+0.000$ & 0/10/0 \\
Bitcoin-OTC$^+$ & $+0.009$ & 0/2/8 & $+0.000$ & 0/9/1 & $+0.000$ & 0/10/0 \\
\texttt{ca-HepTh} & $-0.006$ & 7/2/1 & $-0.006$ & 8/2/0 & $-0.003$ & 5/5/0 \\
\texttt{ca-HepPh} & $-0.002$ & 8/2/0 & $+0.000$ & 2/8/0 & $+0.000$ & 0/10/0 \\
\texttt{ca-AstroPh} & $+0.037$ & 3/0/7 & $+0.009$ & 4/0/6 & $-0.047$ & 10/0/0 \\
\texttt{ca-CondMat} & $-0.019$ & 7/1/2 & $-0.030$ & 9/0/1 & $-0.029$ & 8/2/0 \\
\texttt{email-Enron} & $-0.001$ & 5/0/5 & $-0.013$ & 7/0/3 & $-0.025$ & 9/0/1 \\
\texttt{loc-Brightkite} & $-0.001$ & 5/0/5 & $-0.001$ & 5/1/4 & $-0.001$ & 7/0/3 \\
\texttt{Slashdot+} & $+0.018$ & 0/0/10 & $+0.002$ & 3/0/7 & $-0.001$ & 6/2/2 \\
\texttt{soc-Epinions} & $+0.016$ & 1/0/9 & $+0.004$ & 2/0/8 & $+0.000$ & 4/0/6 \\
\texttt{Epinions+} & $+0.010$ & 1/0/9 & $+0.007$ & 1/0/9 & $+0.000$ & 5/0/5 \\
\texttt{com-DBLP} & $+0.018$ & 1/0/9 & $+0.013$ & 0/0/10 & $-0.004$ & 9/0/1 \\
\texttt{com-Amazon} & $+0.033$ & 1/0/9 & $+0.012$ & 0/0/10 & $+0.003$ & 0/0/10 \\
\texttt{com-Youtube} & $+0.641$ & 2/0/8 & $+0.188$ & 0/0/10 & $+0.009$ & 4/0/6 \\
\bottomrule
\end{tabular}

\end{table}

\paragraph{Dependence on the budget.}  Table~\ref{tab:budget} shows the
limitation of PXMem.  On the two largest graphs KaPoCE is $0.1\%$ to $0.7\%$
better at \SI{60}{s} and \SI{150}{s}, because the initial population (four
Pivot, flipping and annealing runs) takes a large part of a short budget.  On
\texttt{ca-HepTh}, \texttt{ca-AstroPh} and \texttt{ca-CondMat} PXMem is better at
every budget.

\paragraph{Parallel runs and post-processing.}  Theorem~\ref{thm:px} gives a
way to use several cores that heuristics without an exact recombination do not
have.  Table~\ref{tab:parallel} reports four independent PXMem runs of
\SI{300}{s} that ran at the same time on four cores of our 4-core machine,
compared with the best of $k$ of them and with their partition crossover, and
with KaPoCE on one core.  Combining two runs by partition crossover is better
than the best of four runs on four of the five graphs where the runs differ
(all but \texttt{email-Enron}), and never worse than the best of two.  With four
runs PXMem matches or beats KaPoCE on five of the six graphs and is one
disagreement behind on \texttt{com-Amazon}.  Applied to the solutions of
KaPoCE, the crossover improves them on all five graphs where they are not
already equal to ours (by $2$ to $51$ disagreements), which makes it useful as
a post-processing step for any solver.  These numbers compare four cores with
one and complement, but do not replace, the single-core comparison of
Table~\ref{tab:h2h}.

\begin{table}[t]
\centering
\caption{Four independent PXMem runs (\SI{300}{s} each, run in parallel):
best of $k$ runs and partition crossover of $k$ runs (averaged over all
subsets of size $k$), KaPoCE (\SI{300}{s}, one core), and the partition
crossover of the KaPoCE solution with a PXMem solution (averaged over the four
runs).  The best value in each row is bold.}\label{tab:parallel}
\small
\begin{tabular}{lrrrrrrrr}
\toprule
 & 1 run & \multicolumn{2}{c}{2 runs} & \multicolumn{2}{c}{4 runs} & & \multicolumn{2}{c}{PX with KaPoCE} \\
\cmidrule(lr){3-4}\cmidrule(lr){5-6}\cmidrule(lr){8-9}
graph & & best & PX & best & PX & KaPoCE & 1 run & 4 runs \\
\midrule
Bitcoin-OTC$^+$ & \textbf{\num{16462}} & \textbf{\num{16462}} & \textbf{\num{16462}} & \textbf{\num{16462}} & \textbf{\num{16462}} & \textbf{\num{16462}} & \textbf{\num{16462}} & \textbf{\num{16462}} \\
\texttt{ca-AstroPh} & \num{125544.5} & \num{125542.8} & \num{125541.7} & \num{125542} & \textbf{\num{125540}} & \num{125567} & \num{125542.8} & \textbf{\num{125540}} \\
\texttt{email-Enron} & \num{150453.8} & \num{150445.8} & \num{150444.5} & \num{150433} & \textbf{\num{150432}} & \num{150438} & \num{150435.8} & \textbf{\num{150432}} \\
\texttt{loc-Brightkite} & \num{174706.2} & \num{174704.5} & \num{174699.7} & \num{174703} & \num{174697} & \num{174699} & \num{174696} & \textbf{\num{174695}} \\
\texttt{com-DBLP} & \num{606773.8} & \num{606761.2} & \num{606735.3} & \num{606754} & \num{606720} & \num{606779} & \num{606728.2} & \textbf{\num{606719}} \\
\texttt{com-Amazon} & \num{612374.5} & \num{612361.7} & \num{612330.3} & \num{612350} & \num{612317} & \num{612316} & \num{612309.2} & \textbf{\num{612308}} \\
\bottomrule
\end{tabular}

\end{table}

\paragraph{Memory.}  PXMem needs more memory than KaPoCE.  The peak resident
set sizes measured in the runs of Table~\ref{tab:h2h} are \SI{360}{MB} to
\SI{370}{MB} against \SI{180}{MB} to \SI{190}{MB} on \texttt{com-Amazon} and
\texttt{com-DBLP}, and \SI{970}{MB} against \SI{730}{MB} on \texttt{com-Youtube}.
